% Figure: propulsive efficiency of an air-breathing jet against the ratio of
% jet velocity to flight velocity. eta_p = 2/(1 + v_j/v_0). The curve falls
% monotonically as the jet is thrown faster than the aircraft flies, which is
% the quantitative reason high-bypass turbofans (small v_j/v_0) beat turbojets
% (large v_j/v_0) at subsonic cruise.
\begin{figure}[htbp]
\centering
\begin{tikzpicture}
\begin{axis}[
  width=11.5cm, height=7.5cm,
  xlabel={jet-to-flight velocity ratio $v_j/v_0$},
  ylabel={propulsive efficiency $\eta_p$},
  xmin=1, xmax=6, ymin=0, ymax=1.05,
  axis lines=left,
  tick label style={font=\small}, label style={font=\small},
  legend pos=north east, legend cell align=left, legend style={font=\scriptsize},
  clip=false,
]
% eta_p = 2/(1 + r), r = v_j/v_0
\addplot[engblue, very thick, samples=120, domain=1:6] {2/(1 + x)};
\addlegendentry{$\eta_p = 2/(1 + v_j/v_0)$}
% marked operating regions
\addplot[engred, only marks, mark=*, mark size=1.8pt] coordinates {(1.25,0.889)};
\node[engred, font=\scriptsize, anchor=west] at (axis cs:1.32,0.889)
  {high-bypass turbofan};
\addplot[engamber, only marks, mark=*, mark size=1.8pt] coordinates {(2.5,0.571)};
\node[engamber, font=\scriptsize, anchor=west] at (axis cs:2.6,0.571)
  {low-bypass turbofan};
\addplot[engorange, only marks, mark=*, mark size=1.8pt] coordinates {(4.5,0.364)};
\node[engorange, font=\scriptsize, anchor=west] at (axis cs:4.6,0.400)
  {turbojet};
\end{axis}
\end{tikzpicture}
\caption[Propulsive efficiency versus jet-to-flight velocity ratio]{Propulsive
efficiency of an air-breathing engine against the ratio of exhaust jet velocity
to flight velocity. The Froude expression $\eta_p = 2/(1 + v_j/v_0)$ falls
steadily as the jet is thrown faster relative to the aircraft, because the
kinetic energy left in the wake is wasted work. A high-bypass turbofan produces
its thrust by giving a large mass of air a small velocity increment, so it sits
near the left edge where the propulsive efficiency is high; a bare turbojet
throws a small mass fast and pays the penalty on the right. The thermal
efficiency of the core Brayton cycle multiplies this propulsive term to give the
overall efficiency, which is why the two levers, a hot efficient core and a
gently accelerated bypass stream, are pursued together.}
\label{fig:vol04ch03:propulsive-efficiency}
\end{figure}
